\section{Introduction}
\IEEEPARstart{F}{ree-floating} space manipulators are the enabling technology for on-orbit servicing, debris removal, and in-space assembly~\cite{liu2021review,yoshida2003etsvii}. Before capture, the base is not rigidly supported: joint motion produces a reaction motion of the base through momentum conservation, so the end-effector is coupled to the entire spacecraft--manipulator system~\cite{dubowsky1993,papadopoulos1991nature}. This coupling, together with parameter uncertainty, disturbances, and actuator limits, makes trajectory tracking substantially harder than for a fixed-base arm, and the difficulty is aggravated when the settling time is a task requirement rather than a by-product of the gains.

\emph{Related work.} Three lines are relevant. First, the coupled kinematics is described by the generalized Jacobian~\cite{umetani1989}, whose configuration-dependent singularity loci are now well understood~\cite{papadopoulos1993,nanos2015avoiding}, and the base motion can be eliminated to obtain an equivalent joint-space model~\cite{parlaktuna2004}. Second, robust tracking has been addressed by disturbance observers~\cite{chu2008dob}, adaptive Jacobian schemes~\cite{wang2012adaptivejacobian}, time-delay estimation~\cite{wang2025tdelog}, and terminal sliding modes with finite-time observers~\cite{shao2021ntsm,dou2023fractional,dou2025fractional}. Third, the convergence time itself has become a design objective: finite-time~\cite{jia2020finite}, fixed-time~\cite{polyakov2012,hu2024constrainedfixed}, and predefined- or prescribed-time designs~\cite{xu2024se3predefined,zou2025practical,yan} specify it explicitly. Most of the last group relies on fractional powers, terminal nonsmooth terms, or gains that grow without bound near the terminal instant, and the disturbance observers they use offer no handle on the observation time~\cite{yan}. Saturation and performance envelopes add further conservatism when several constraints coexist~\cite{zhou2025saturated}.

\emph{Delayed feedback.} Smooth periodic delayed feedback (PDF) is an alternative that has no terminal singularity: for a controllable linear time-delay system, a periodic delayed gain annihilates the one-period state-transition operator, which yields fixed-time stabilisation with a continuous input~\cite{pyragas1992,zhouFixedtimeStabilizationLinear2022}. Extending this mechanism to a free-floating manipulator requires three ingredients that are not available in~\cite{zhouFixedtimeStabilizationLinear2022}: an exact auxiliary input channel that absorbs the nonlinear and momentum-coupled terms; a disturbance interface whose observation time is designable, for which a prescribed-time observer is the natural candidate~\cite{jiangPrescribedtimeDisturbanceObserver2024}; and a quantification of how the residual observation error, saturation, and sampling perturb the one-period map.

\emph{Contributions.} This correspondence
\begin{enumerate}
\item reduces the free-floating manipulator to an exact controllable error channel using the Schur complement of the spatial inertia matrix and the momentum-consistent bias forces, and designs a smooth PDF law on that channel;
\item states the one-period nullification theorem for the reduced nonlinear plant, together with a practical bound in which the observation residual, the actuator limit, and the discretisation appear explicitly;
\item reports a design rule for the Gramian gain: because the equivalent inertia of a long-reach arm spans three decades, normalising the gains by the joint inertia spreads the modal time scales and makes the Gramian numerically singular, whereas a uniform gain keeps the Gramian block diagonal and well conditioned;
\item evaluates the design on a complete three-dimensional model of a seven-degree-of-freedom manipulator with published parameters and $\pm20~\mathrm{N\,m}$ actuators, reporting settling time, peak torque, saturation duty cycle, and base reaction motion.
\end{enumerate}

\emph{Organization.} Section~\ref{sec:model} states the reduced model, Section~\ref{sec:control} the observer and the PDF law, Section~\ref{sec:stability} the stability and practical bound, Section~\ref{sec:simulation} the numerical example, and Section~\ref{sec:conclusion} the conclusions. Long derivations are collected in the Appendix.
